
\documentclass{article}
\usepackage{colortbl}
\usepackage{makecell}
\usepackage{multirow}
\usepackage{supertabular}

\begin{document}

\newcounter{utterance}

\centering \large Interaction Transcript for game `dond', experiment `coop\_en', episode 4 with Qwen3.5{-}27B.
\vspace{24pt}

{ \footnotesize  \setcounter{utterance}{1}
\setlength{\tabcolsep}{0pt}
\begin{supertabular}{c@{$\;$}|p{.15\linewidth}@{}p{.15\linewidth}p{.15\linewidth}p{.15\linewidth}p{.15\linewidth}p{.15\linewidth}}
    \# & $\;$A & \multicolumn{4}{c}{Game Master} & $\;\:$B\\
    \hline

    \theutterance \stepcounter{utterance}  
    & & \multicolumn{4}{p{0.6\linewidth}}{
        \cellcolor[rgb]{0.9,0.9,0.9}{
            \makecell[{{p{\linewidth}}}]{
                \texttt{\tiny{[P1$\langle$GM]}}
                \texttt{You are playing a cooperative negotiation game in which you have to agree on how to divide a set of items among you and another player.} \\
\\ 
\texttt{Rules:} \\
\texttt{(a) You and the other player are given a set of items. Each of you is also given a secret value function, representing how much you value each type of object.} \\
\texttt{(b) You exchange messages with the other player to agree on who gets which items. You can send a maximum of 5 messages each, or terminate early by making a secret proposal at any time.} \\
\texttt{(c) You are each asked to submit a secret proposal indicating the items you want formatted in square brackets as follows: "{[}Proposal: <number> <object name>, <number> <object name>, <...>{]}"} \\
\texttt{(d) If your proposals are complementary, i.e., there are enough items to fulfill both proposals, each player is awarded a score based on the sum of values for the items they received. Otherwise, both of you get zero points.} \\
\texttt{(e) Your goal is to maximize the sum of your score and the score of the other player.} \\
\\ 
\texttt{Let us start.} \\
\\ 
\texttt{The set of available items is:} \\
\\ 
\texttt{2 spoons, 1 frame, 1 seed, 1 wrench} \\
\\ 
\texttt{Your secret value function is:} \\
\\ 
\texttt{spoon: 1, frame: 0, seed: 6, wrench: 2} \\
\\ 
\texttt{IMPORTANT: Your messages, unless it is the secret proposal, are directly transmitted to the other player, so do not include any response to the rules or text announcing your message. To make a secret proposal, use the indicated format. Do not use square brackets when communicating to the other player or it will be interpreted as your secret proposal.} \\
            }
        }
    }
    & \\ \\

    \theutterance \stepcounter{utterance}  
    & \multicolumn{4}{p{0.6\linewidth}}{
        \cellcolor[rgb]{0.9,0.9,0.9}{
            \makecell[{{p{\linewidth}}}]{
                \texttt{\tiny{[P1$\rangle$GM]}}
                \texttt{Hi! Let's figure out the best way to split these items. I value the seed the most (6 points), followed by the wrench (2 points), and the spoons are worth 1 each. I don't really care about the frame (0 points).} \\
\\ 
\texttt{What are your values for these items? Knowing that will help us make sure we both get the things we want most.} \\
            }
        }
    }
    & & \\ \\

    \theutterance \stepcounter{utterance}  
    & & & \multicolumn{2}{p{0.3\linewidth}}{
        \cellcolor[rgb]{0.9,0.9,0.9}{
            \makecell[{{p{\linewidth}}}]{
                \texttt{\tiny{[GM$|$GM]}}
                \texttt{continue} \\
            }
        }
    }
    & & \\ \\

    \theutterance \stepcounter{utterance}  
    & & & \multicolumn{2}{p{0.3\linewidth}}{
        \cellcolor[rgb]{0.9,0.9,0.9}{
            \makecell[{{p{\linewidth}}}]{
                \texttt{\tiny{[GM$|$GM]}}
                \texttt{Hi! Let's figure out the best way to split these items. I value the seed the most (6 points), followed by the wrench (2 points), and the spoons are worth 1 each. I don't really care about the frame (0 points).} \\
\\ 
\texttt{What are your values for these items? Knowing that will help us make sure we both get the things we want most.} \\
            }
        }
    }
    & & \\ \\

    \theutterance \stepcounter{utterance}  
    & & \multicolumn{4}{p{0.6\linewidth}}{
        \cellcolor[rgb]{0.9,0.9,0.9}{
            \makecell[{{p{\linewidth}}}]{
                \texttt{\tiny{[GM$\rangle$P2]}}
                \texttt{You are playing a cooperative negotiation game in which you have to agree on how to divide a set of items among you and another player.} \\
\\ 
\texttt{Rules:} \\
\texttt{(a) You and the other player are given a set of items. Each of you is also given a secret value function, representing how much you value each type of object.} \\
\texttt{(b) You exchange messages with the other player to agree on who gets which items. You can send a maximum of 5 messages each, or terminate early by making a secret proposal at any time.} \\
\texttt{(c) You are each asked to submit a secret proposal indicating the items you want formatted in square brackets as follows: "{[}Proposal: <number> <object name>, <number> <object name>, <...>{]}"} \\
\texttt{(d) If your proposals are complementary, i.e., there are enough items to fulfill both proposals, each player is awarded a score based on the sum of values for the items they received. Otherwise, both of you get zero points.} \\
\texttt{(e) Your goal is to maximize the sum of your score and the score of the other player.} \\
\\ 
\texttt{Let us start.} \\
\\ 
\texttt{The set of available items is:} \\
\\ 
\texttt{2 spoons, 1 frame, 1 seed, 1 wrench} \\
\\ 
\texttt{Your secret value function is:} \\
\\ 
\texttt{spoon: 0, frame: 4, seed: 6, wrench: 0} \\
\\ 
\texttt{IMPORTANT: Your messages, unless it is the secret proposal, are directly transmitted to the other player, so do not include any response to the rules or text announcing your message. To make a secret proposal, use the indicated format. Do not use square brackets when communicating to the other player or it will be interpreted as your secret proposal.} \\
\\ 
\\ 
\texttt{Hi! Let's figure out the best way to split these items. I value the seed the most (6 points), followed by the wrench (2 points), and the spoons are worth 1 each. I don't really care about the frame (0 points).} \\
\\ 
\texttt{What are your values for these items? Knowing that will help us make sure we both get the things we want most.} \\
            }
        }
    }
    & \\ \\

    \theutterance \stepcounter{utterance}  
    & & & \multicolumn{4}{p{0.6\linewidth}}{
        \cellcolor[rgb]{0.9,0.9,0.9}{
            \makecell[{{p{\linewidth}}}]{
                \texttt{\tiny{[GM$\langle$P2]}}
                \texttt{Hi! Thanks for sharing your values. Here's how I value the items:} \\
\\ 
\texttt{{-} **Seed**: 6 points (highest priority for me)} \\
\texttt{{-} **Frame**: 4 points} \\
\texttt{{-} **Spoons**: 0 points each} \\
\texttt{{-} **Wrench**: 0 points} \\
\\ 
\texttt{Since you value the seed most (6) and I also value it most (6), we need to decide who gets it. You also value the wrench (2) and spoons (1 each), while I don't care about those. I value the frame (4), which you don't care about.} \\
\\ 
\texttt{To maximize our combined score:} \\
\texttt{{-} If I take the **seed** (6) and **frame** (4), my total is 10.} \\
\texttt{{-} You take the **wrench** (2) and **2 spoons** (2), your total is 4.} \\
\texttt{{-} Combined: 14 points.} \\
\\ 
\texttt{Alternatively:} \\
\texttt{{-} If you take the **seed** (6) and **wrench** (2) and **2 spoons** (2), your total is 10.} \\
\texttt{{-} I take the **frame** (4), my total is 4.} \\
\texttt{{-} Combined: 14 points.} \\
\\ 
\texttt{Both options give the same total, but since you mentioned you value the seed most, maybe you should take it? Or we can split based on who values the frame more (me). Let me know your thoughts!} \\
            }
        }
    }
     \\ \\

    \theutterance \stepcounter{utterance}  
    & & & \multicolumn{2}{p{0.3\linewidth}}{
        \cellcolor[rgb]{0.9,0.9,0.9}{
            \makecell[{{p{\linewidth}}}]{
                \texttt{\tiny{[GM$|$GM]}}
                \texttt{continue} \\
            }
        }
    }
    & & \\ \\

    \theutterance \stepcounter{utterance}  
    & & & \multicolumn{2}{p{0.3\linewidth}}{
        \cellcolor[rgb]{0.9,0.9,0.9}{
            \makecell[{{p{\linewidth}}}]{
                \texttt{\tiny{[GM$|$GM]}}
                \texttt{Hi! Thanks for sharing your values. Here's how I value the items:} \\
\\ 
\texttt{{-} **Seed**: 6 points (highest priority for me)} \\
\texttt{{-} **Frame**: 4 points} \\
\texttt{{-} **Spoons**: 0 points each} \\
\texttt{{-} **Wrench**: 0 points} \\
\\ 
\texttt{Since you value the seed most (6) and I also value it most (6), we need to decide who gets it. You also value the wrench (2) and spoons (1 each), while I don't care about those. I value the frame (4), which you don't care about.} \\
\\ 
\texttt{To maximize our combined score:} \\
\texttt{{-} If I take the **seed** (6) and **frame** (4), my total is 10.} \\
\texttt{{-} You take the **wrench** (2) and **2 spoons** (2), your total is 4.} \\
\texttt{{-} Combined: 14 points.} \\
\\ 
\texttt{Alternatively:} \\
\texttt{{-} If you take the **seed** (6) and **wrench** (2) and **2 spoons** (2), your total is 10.} \\
\texttt{{-} I take the **frame** (4), my total is 4.} \\
\texttt{{-} Combined: 14 points.} \\
\\ 
\texttt{Both options give the same total, but since you mentioned you value the seed most, maybe you should take it? Or we can split based on who values the frame more (me). Let me know your thoughts!} \\
            }
        }
    }
    & & \\ \\

    \theutterance \stepcounter{utterance}  
    & & \multicolumn{4}{p{0.6\linewidth}}{
        \cellcolor[rgb]{0.9,0.9,0.9}{
            \makecell[{{p{\linewidth}}}]{
                \texttt{\tiny{[P1$\langle$GM]}}
                \texttt{Hi! Thanks for sharing your values. Here's how I value the items:} \\
\\ 
\texttt{{-} **Seed**: 6 points (highest priority for me)} \\
\texttt{{-} **Frame**: 4 points} \\
\texttt{{-} **Spoons**: 0 points each} \\
\texttt{{-} **Wrench**: 0 points} \\
\\ 
\texttt{Since you value the seed most (6) and I also value it most (6), we need to decide who gets it. You also value the wrench (2) and spoons (1 each), while I don't care about those. I value the frame (4), which you don't care about.} \\
\\ 
\texttt{To maximize our combined score:} \\
\texttt{{-} If I take the **seed** (6) and **frame** (4), my total is 10.} \\
\texttt{{-} You take the **wrench** (2) and **2 spoons** (2), your total is 4.} \\
\texttt{{-} Combined: 14 points.} \\
\\ 
\texttt{Alternatively:} \\
\texttt{{-} If you take the **seed** (6) and **wrench** (2) and **2 spoons** (2), your total is 10.} \\
\texttt{{-} I take the **frame** (4), my total is 4.} \\
\texttt{{-} Combined: 14 points.} \\
\\ 
\texttt{Both options give the same total, but since you mentioned you value the seed most, maybe you should take it? Or we can split based on who values the frame more (me). Let me know your thoughts!} \\
            }
        }
    }
    & \\ \\

    \theutterance \stepcounter{utterance}  
    & \multicolumn{4}{p{0.6\linewidth}}{
        \cellcolor[rgb]{0.9,0.9,0.9}{
            \makecell[{{p{\linewidth}}}]{
                \texttt{\tiny{[P1$\rangle$GM]}}
                \texttt{That's a great analysis! You are absolutely right that both scenarios yield a combined score of 14, which is the maximum possible.} \\
\\ 
\texttt{Since we both value the seed equally (6 points), the deciding factor is the frame. You value the frame at 4 points, while I value it at 0. Therefore, it makes the most sense for you to take the frame.} \\
\\ 
\texttt{To maximize the total utility, let's go with the option where you get the items you value most (Seed and Frame) and I get the items I value (Wrench and Spoons).} \\
\\ 
\texttt{Here is the plan:} \\
\texttt{{-} **You**: 1 seed, 1 frame} \\
\texttt{{-} **Me**: 1 wrench, 2 spoons} \\
\\ 
\texttt{Does this work for you? If so, I will submit the secret proposal for my items.} \\
\\ 
\texttt{{[}Proposal: 1 wrench, 2 spoons{]}} \\
            }
        }
    }
    & & \\ \\

    \theutterance \stepcounter{utterance}  
    & & & \multicolumn{2}{p{0.3\linewidth}}{
        \cellcolor[rgb]{0.9,0.9,0.9}{
            \makecell[{{p{\linewidth}}}]{
                \texttt{\tiny{[GM$|$GM]}}
                \texttt{proposal} \\
            }
        }
    }
    & & \\ \\

    \theutterance \stepcounter{utterance}  
    & & & \multicolumn{2}{p{0.3\linewidth}}{
        \cellcolor[rgb]{0.9,0.9,0.9}{
            \makecell[{{p{\linewidth}}}]{
                \texttt{\tiny{[GM$|$GM]}}
                [2, 0, 0, 1]
            }
        }
    }
    & & \\ \\

    \theutterance \stepcounter{utterance}  
    & & \multicolumn{4}{p{0.6\linewidth}}{
        \cellcolor[rgb]{0.9,0.9,0.9}{
            \makecell[{{p{\linewidth}}}]{
                \texttt{\tiny{[GM$\rangle$P2]}}
                \texttt{The other player has submitted a secret proposal. It is now your turn to submit a proposal of your own.} \\
            }
        }
    }
    & \\ \\

    \theutterance \stepcounter{utterance}  
    & & & \multicolumn{4}{p{0.6\linewidth}}{
        \cellcolor[rgb]{0.9,0.9,0.9}{
            \makecell[{{p{\linewidth}}}]{
                \texttt{\tiny{[GM$\langle$P2]}}
                \texttt{{[}Proposal: 1 frame, 1 seed{]}} \\
            }
        }
    }
     \\ \\

    \theutterance \stepcounter{utterance}  
    & & & \multicolumn{2}{p{0.3\linewidth}}{
        \cellcolor[rgb]{0.9,0.9,0.9}{
            \makecell[{{p{\linewidth}}}]{
                \texttt{\tiny{[GM$|$GM]}}
                \texttt{proposal} \\
            }
        }
    }
    & & \\ \\

    \theutterance \stepcounter{utterance}  
    & & & \multicolumn{2}{p{0.3\linewidth}}{
        \cellcolor[rgb]{0.9,0.9,0.9}{
            \makecell[{{p{\linewidth}}}]{
                \texttt{\tiny{[GM$|$GM]}}
                [0, 1, 1, 0]
            }
        }
    }
    & & \\ \\

    \theutterance \stepcounter{utterance}  
    & & & \multicolumn{2}{p{0.3\linewidth}}{
        \cellcolor[rgb]{0.9,0.9,0.9}{
            \makecell[{{p{\linewidth}}}]{
                \texttt{\tiny{[GM$|$GM]}}
                [[2, 0, 0, 1], [0, 1, 1, 0]]
            }
        }
    }
    & & \\ \\

\end{supertabular}
}

\end{document}
